\chapter{Fourier Series and Fourier Transform}
\renewcommand{\currentchapterpath}{chapters/ch05/}

\begin{flushright}
\textit{“八音克谐，无相夺伦。” ——《尚书·舜典》}
\footnote{From the ``Canon of Shun'' in the \textit{Book of Documents}: ``The eight classes of musical sounds achieve harmony, with none displacing the proper order of another.'' The line praises harmony formed from distinct components. Fourier analysis similarly represents a complicated signal as a superposition of elementary oscillations, each with its own frequency and amplitude; their ordered combination reconstructs the original behavior.}
\end{flushright}

Fourier series and Fourier transforms resolve functions into elementary oscillations. They occur throughout signal processing, acoustics, optics, probability, statistics, and the solution of differential equations. This chapter develops Fourier series for periodic functions, passes to the Fourier transform for functions on the real line, and then introduces the Laplace transform as a closely related tool for initial-value problems.

\chapterinput{fourier_series}
\chapterinput{fourier_transform}
\chapterinput{delta_function}
\chapterinput{laplace_transform}
\chapterinput{applications}
\chapterinput{exercises}